%% =============================================================================
%% SEC 3 - IS IT ALPHA? ~830 words + Table 2.
%% SOURCE: paper1_sec4_benchmarks.tex + paper1_sec5_factor_selection.tex.
%% v2 fixes: (i) selector count -- knockoffs is ONE of the five alternatives,
%% as in the paper; (ii) regime multiplier restated as 3-7x across the NINE
%% benchmark cells (iTraxx best-subset is flatter, so "every control set" was
%% wrong); (iii) added sub-period + rolling-alpha stability sentence and the
%% recurring-exposures / double-selection line from Section III.C.
%% =============================================================================

The question is whether the returns of Table~\ref{tab:performance} are
compensation for systematic risk. We answer it with an escalating battery: if
the alpha is an artifact of an omitted exposure, some layer should absorb it.
% SOURCE: sec4 opening logic.

The first three layers are the canonical benchmarks of the fixed-income
arbitrage and active-management literatures---the ten-factor model of
\citet{duarte2007risk}, the seven-factor model of \citet{fung2004hedge}, and
the eight risk premia of \citet{brooks2020active}---estimated on euro
counterparts of every factor (U.S.-factor versions, with the same conclusions,
in Appendix~\ref{app:benchmarks}). The fourth layer replaces specified factors
with the first eight principal components of a candidate universe of 75
tradable risk factors spanning credit, liquidity, equity, volatility, and
rates, with $K=8$ selected by the information criteria of
\citet{bai2002determining}. The fifth is best-subset selection \citep{bertsimas2016best}: for each
strategy, an $\ell_0$-constrained search picks whichever eight of the
candidates explain its returns best---the hardest test we can pose, because
the data are free to construct the most absorbing hedge. Formally,
\begin{equation}
(\hat\alpha_i,\hat{\boldsymbol\beta}_i)=\arg\min_{\alpha,\,\boldsymbol\beta}\;
\sum_{t}\bigl(r_{i,t}-\alpha-\boldsymbol\beta^{\prime}\mathbf{F}_t\bigr)^{2}
\quad\text{s.t.}\quad \lVert\boldsymbol\beta\rVert_0\le k,\qquad k=8,
\end{equation}
where $\lVert\boldsymbol\beta\rVert_0$ counts the nonzero coefficients.
% SOURCE: sec4 A + sec5 A-C, compressed.

Every layer estimates the spanning regression
\begin{equation}
r_{i,t} = \alpha_i + \sum_{j}\beta_{i,j}\,F_{j,t} + \varepsilon_{i,t},
\end{equation}
where $r_{i,t}$ is the monthly net excess return of strategy $i$, $F_{j,t}$
the layer's factor set, and the null of interest is $\alpha_i = 0$.
Table~\ref{tab:alpha_layers} reports the outcome. Moving down the rows,
explanatory power rises from near zero to adjusted $R^2$ of 21--31\%, while
the alpha barely moves: 1.4--2.2\% per year for BTP~Italia, 4.0--4.8\% for the
CDS--bond basis, 1.7--2.9\% for the iTraxx skew, significant at the 1\% level
in every cell. Panel~B repeats the exercise in real time: at each month the
hedge is re-estimated on a past-only expanding window
\citep{welch2008comprehensive}, so the surviving alpha is an abnormal return
an investor could actually have earned. It survives everywhere, and for the
best-subset hedge every out-of-sample $t$-statistic clears the 3.0 hurdle that
\citet{harvey2016and} propose for discoveries in multiple-testing
environments. The estimate is also stable through time: it is significant in
every sub-period we examine, and 36-month rolling alphas are positive across
essentially the entire sample under all three benchmarks
(Appendix~\ref{app:benchmarks}).
% SOURCE: sec4 B closing + Figure 4 discussion.

Three details sharpen the verdict. First, the three selected factor sets are
pairwise disjoint---no factor appears in more than one strategy---so no single
common exposure is responsible; the economically telling loading is the iTraxx
skew's on the traded intermediary-capital factor of \citet{he2017intermediary}:
the strategy pays off when dealer equity underperforms. Second, the result is
invariant to the selector. Five alternative selectors span the methodological
space---a best subset constrained to one factor per risk category, the
adaptive elastic net \citep{zou2009adaptive}, the relaxed lasso
\citep{meinshausen2007relaxed}, a machine-learning screen based on random
forests \citep{breiman2001random},
and model-X knockoffs \citep{candes2018panning} with the false discovery rate
controlled at 10\%---and all leave the alpha intact under honest split-sample
inference, with factors chosen on the first half of the sample and alpha
estimated on the second. The knockoffs are the sharpest statement of the
two-sided message. Each candidate $j$ receives a statistic $W_j$ contrasting
its explanatory strength against a synthetic knockoff copy, and the selection
threshold
\begin{equation}
\tau_q=\min\Bigl\{t>0:\
\frac{1+\#\{j:\,W_j\le -t\}}{\max\bigl(1,\,\#\{j:\,W_j\ge t\}\bigr)}\le q
\Bigr\},\qquad q=0.10,
\end{equation}
controls the false discovery rate at $q$. In our data no threshold satisfies
the bound: the procedure selects nothing, so the data neither identify a
stable factor structure nor produce one that absorbs the premium
(Appendix~\ref{app:selection}). The exposures that recur across selectors are
the consensus ones---fixed-income defensive for BTP~Italia, emerging-market
debt for the CDS--bond basis, the intermediary block for the iTraxx skew---and
the post-double-selection control set of \citet{belloni2014inference} is a
strict subset of exactly these. Third, the alpha is state-dependent. The conditional framework of
\citet{christopherson1998conditioning} lets the alpha and the loadings vary
with a lagged, publicly observable state variable,
\begin{equation}
r_{i,t} = \alpha_0 + \alpha_1 z_{t-1}
        + (\boldsymbol\beta + \boldsymbol\delta\, z_{t-1})^{\prime}\mathbf{F}_t
        + \varepsilon_{i,t},
\end{equation}
where $z_{t-1}$ is the standardized iTraxx Main 5-year spread. Estimated
across the frameworks, a one-standard-deviation rise in $z$ adds 1.7--3.2
percentage points of annual alpha to the CDS--bond basis and 0.9--1.5 to the
iTraxx skew; split discretely by regime, the high-stress alpha is three to
seven times the calm-market alpha---a gap of 2.3 to 5.1 percentage points per
year---across the nine strategy--benchmark
combinations, and rises monotonically with stress under the data-driven
controls as well (Appendices~\ref{app:benchmarks}--\ref{app:selection}).
% SOURCE: sec5 C-D + sec4 B; numbers from Tables VI, X, XI, XIII.

An alpha that no factor spans, that survives in real time, and that
concentrates in funding stress is not a risk premium in disguise; it is the
fingerprint of an arbitrage that capital fails to close precisely when closing
it is most profitable. The next section asks why.
% ALEX-WRITE: check the bridge sentence lands in your voice.

%% TABLE 2 - ALPHA ACROSS LAYERS, IN AND OUT OF SAMPLE (compact).
%% Panel A = paper Table XII (verified cell-by-cell against the compiled paper).
%% Panel B = paper Table VI Panel C + Table IX expanding-window OOS (verified).
%% Stars from the source tables; Panel A all ***; Panel B: BTP-Duarte **.
\begin{table}[t]
\centering
\caption{Annualized alpha (\% p.a.) across five layers of factor controls.
Panel A: full-sample spanning regressions (Newey--West $t$-statistics; adjusted
$R^2$ in brackets). Panel B: out-of-sample alpha with hedge loadings estimated
on past-only expanding windows (60-month burn-in). $^{***}$, $^{**}$, $^{*}$:
1\%, 5\%, 10\%. All best-subset out-of-sample $t$-statistics exceed the
\citet{harvey2016and} threshold of 3.0.
\label{tab:alpha_layers}}
\small
\begin{tabular}{lccc}
\toprule
Control layer & BTP Italia & CDS--Bond & iTraxx Skew \\
\midrule
\multicolumn{4}{l}{\emph{Panel A: In-sample alpha [adj.\ $R^2$]}}\\
Duarte et al.\ (2007), $k{=}10$   & $+1.6^{***}$ [0.11] & $+4.2^{***}$ [0.15] & $+2.2^{***}$ [0.06] \\
Brooks et al.\ (2020), $k{=}8$    & $+2.1^{***}$ [0.08] & $+4.2^{***}$ [0.12] & $+1.7^{***}$ [0.04] \\
Fung--Hsieh (2004), $k{=}7$       & $+1.8^{***}$ [0.10] & $+4.0^{***}$ [0.11] & $+1.8^{***}$ [0.00] \\
Principal components, $K{=}8$     & $+1.4^{***}$ [0.09] & $+4.8^{***}$ [0.11] & $+2.2^{***}$ [0.13] \\
Best subset ($\ell_0$), $k{=}8$   & $+2.2^{***}$ [0.31] & $+4.8^{***}$ [0.21] & $+2.9^{***}$ [0.25] \\
\midrule
\multicolumn{4}{l}{\emph{Panel B: Out-of-sample alpha (expanding window)}}\\
Duarte et al.\ (2007)             & $+1.29^{**}$  & $+3.48^{***}$ & $+2.27^{***}$ \\
Brooks et al.\ (2020)             & $+2.10^{***}$ & $+3.86^{***}$ & $+1.77^{***}$ \\
Fung--Hsieh (2004)                & $+1.69^{***}$ & $+3.45^{***}$ & $+2.02^{***}$ \\
Principal components              & $+1.61^{***}$ & $+4.55^{***}$ & $+1.85^{***}$ \\
Best subset (frozen set)          & $+2.05^{***}$ & $+3.90^{***}$ & $+2.50^{***}$ \\
\bottomrule
\end{tabular}
\end{table}

